\chapter{Conclusion}
\label{ch:conclusion}

\section{Summary of Contributions}

This thesis examined the configuration parameters that govern the cost of a nonuniform fast Fourier transform in FINUFFT, and measured how each of them acts on the individual steps of the algorithm rather than on the transform as a whole. Stage-level timing first established that the spreading and interpolation phase and the FFT together account for essentially all of the work in types 1 and 2, which is where every parameter studied here has its effect. A Gaussian-process search over the joint space of six parameters then ranked them by their contribution to the variance in execution time, and the chapters that followed took the highest-ranked of them apart mechanism by mechanism: the upsampling factor $\sigma$ in Chapter~\ref{ch:sigma_selection}, and the maximum subproblem size and the sort bin dimensions in Chapter~\ref{ch:parameter_tuning}.

\subsection{Conclusion}

The common finding across all of them is that none of these parameters has a monotone effect. Each one accelerates one step of the algorithm and slows another, so the cost as a function of the parameter is a sum of two opposing terms and its minimum lies at the point where they balance. The upsampling factor sets the width of the kernel and the size of the fine grid at the same time: raising it narrows the kernel and makes spreading cheaper, while enlarging the grid the FFT has to transform. The maximum subproblem size trades the redundant write traffic of many small subgrids against the load imbalance of few large ones, and shifts work between the gathering and adding steps in opposite directions. The sort bin dimensions trade the cost of the bin sort itself, which is governed by how many bins the counter array has to hold, against the size of the subgrids that the resulting point ordering produces, which the adding step must write out in full; and they do so subject to a third constraint, that the region a bin covers must stay small enough for the spreading step to write into it without leaving the first-level cache.

The position of these balance points is not a property of the library but of the problem and the machine it runs on. Point density decides how many points a subproblem draws from a given region of the grid, and therefore how the cost of writing a subgrid compares with the cost of filling it; the requested tolerance fixes the kernel width, and with it both the spreading work per point and the padding around every subgrid; and the cache hierarchy sets hard thresholds that several of the parameters cross, the sort counter array against the first-level data cache and the region covered by a bin against the first and second levels. Chapters~\ref{ch:sigma_selection} and~\ref{ch:parameter_tuning} give closed-form models for the quantities these thresholds apply to -- the cost of the two phases $\sigma$ trades against each other, and the average subgrid size a given subproblem size and bin shape produce -- and validate them against measurements taken from the spreader itself.

\section{Limitations}

The measurements throughout were taken on a single class of machine, and the thresholds that turn out to matter most are expressed in terms of that machine's cache sizes. The models are written in terms of those sizes rather than in terms of absolute numbers, so they transfer in principle, but this was not verified on hardware with a different cache hierarchy or a different core count. The parameter search of Chapter~\ref{ch:parameter_importance} likewise covers a single workload; its ranking of the parameters is specific to that workload, and the analysis chapters show directly that a parameter which is insensitive at one density and thread count is decisive at another. Finally, the sort bin dimensions, are not exposed by the library at all, and external interface had to be altered to modify them.

\section{Future Work}

The natural continuation of this work is to turn the balance points into predictions. Each of the parameters studied here has a cost model with the right shape -- two opposing terms whose sum has an interior minimum -- but the constants that place that minimum are properties of the machine: the cost of a cache miss at each level, the effective bandwidth available to a thread, the overhead of a subproblem. An empirical model that calibrates these constants once per machine, from a short set of measurements at known configurations, could then predict the minimising parameter values for an arbitrary input from its density, dimension, tolerance and thread count, without a search. The evidence gathered here suggests this is achievable: the subgrid-size model already predicts the measured geometry to within a percent across a wide range of configurations, the cost of the adding step is proportional to the number of cells written with a constant that does not vary, and the transitions in the spreading and sorting steps occur where the corresponding working set crosses a known cache boundary. What remains is to fit the machine-dependent constants that connect these quantities to time, and to validate the resulting predictions against a search of the kind performed in Chapter~\ref{ch:parameter_importance} across a range of inputs and machines.

A second, and prerequisite, line of work is a benchmark suite of transforms broad enough to measure whether such a model actually helps. The parameters studied here are chosen by heuristics, and a heuristic that improves the common case can quietly make some other case much worse: the search of Chapter~\ref{ch:parameter_importance} covered a single workload, and the analysis chapters show repeatedly that a parameter which barely matters at one density, dimension or thread count is decisive at another. A suite spanning those axes -- the three dimensions, densities from far below one point per cell to far above, tolerances across the usable range, single-threaded through fully-threaded, and both precisions -- run against the library defaults would give a baseline against which any proposed heuristic can be scored, reporting not only the mean improvement but the worst regression it causes. That second number is the one that matters for a library: a change that is faster on average and several times slower somewhere is not an improvement, and without a suite of this kind such a regression is only discovered by whoever runs into it.
